% !TEX program = xelatex
% Example: Scientific poster with the beamerposter package
% Uses an A0 portrait layout with a three-column block arrangement.
\documentclass[final]{beamer}
\usepackage[size=a0, scale=1.4, orientation=portrait]{beamerposter}
\usepackage{tikz}
\usetheme{default}
\setbeamertemplate{navigation symbols}{}

% Block styling for poster boxes
\setbeamertemplate{block begin}{
  \begin{beamercolorbox}[ht=2.8ex, dp=1.2ex, leftskip=0.5cm, rightskip=0.5cm]{block title}
    \usebeamerfont*{block title}\large\insertblocktitle
  \end{beamercolorbox}
  \begin{beamercolorbox}[dp=1ex, leftskip=0.5cm, rightskip=0.5cm]{block body}
    \usebeamerfont*{block body}\large
}
\setbeamertemplate{block end}{
  \end{beamercolorbox}\vspace{0.5ex}
}

\title{Efficient Algorithms for Large-Scale Graph Analysis}
\author{Dr.~Alan Turing \quad Dr.~Grace Hopper}
\institute{Institute of Computing, University of Example}
\date{2025}

\begin{document}
\begin{frame}[t]

% Title banner
\begin{beamercolorbox}[wd=\paperwidth, sep=1cm, center]{title}
  {\veryhuge\textbf{Efficient Algorithms for Large-Scale Graph Analysis}}\\[0.5cm]
  {\large Dr.~Alan Turing \quad Dr.~Grace Hopper}\\[0.2cm]
  {\normalsize Institute of Computing, University of Example \quad\textbar\quad 2025}
\end{beamercolorbox}
\vspace{1cm}

\begin{columns}[t]

\begin{column}{0.31\paperwidth}
  \begin{block}{Introduction}
    Graph-structured data appears in social networks, biology,
    and transportation. We study scalable algorithms for
    community detection and shortest-path queries on graphs
    with millions of nodes.
  \end{block}

  \begin{block}{Motivation}
    \begin{itemize}
      \item Real-world graphs are large and sparse.
      \item Exact solutions are often too slow.
      \item Approximate methods trade accuracy for speed.
    \end{itemize}
  \end{block}
\end{column}

\begin{column}{0.31\paperwidth}
  \begin{block}{Method}
    We propose a two-phase approach:
    \begin{enumerate}
      \item Coarsen the graph by merging similar nodes.
      \item Solve on the coarse graph, then refine.
    \end{enumerate}
    The refinement step uses a local search bounded by
    a radius parameter $r$.
  \end{block}

  \begin{block}{Complexity}
    The running time is
    \[
      O\!\left(\frac{n \log n}{r}\right),
    \]
    compared to $O(n^2)$ for the naive baseline.
  \end{block}
\end{column}

\begin{column}{0.31\paperwidth}
  \begin{block}{Results}
    \begin{itemize}
      \item 12$\times$ speedup on road networks.
      \item Within 3\% of the exact solution on average.
      \item Scales to 10M nodes on a single machine.
    \end{itemize}
  \end{block}

  \begin{block}{Conclusion}
    The coarsen-and-refine strategy offers a favorable
    accuracy--speed trade-off for large graph problems.
    Future work will extend the method to dynamic graphs.
  \end{block}
\end{column}

\end{columns}

\end{frame}
\end{document}
